\documentclass{report}
\usepackage{cmt}

\title{\huge{Algorithmische Mathematik I \\ Übungszettel 5}}
\author{\huge{Amy, Rojan}}
\date{2026-11-19}

\begin{document}

\maketitle

\nt{
  Den Übungszettel gibt es
  \href{run:./pset_2026-11-12.pdf}{hier}.
}

\qs{}{
  Zeigen Sie, dass man mithilfe von Karatsubas Algorithmus
  zwei natürliche Zahlen mit $k$- bzw. $l$-stelliger
  Binärdarstellung, wobei $k\le l$ gilt, in
  \[
    O\left(k^{\log_2 3-1}\cdot l\right)
  \]
  Zeit multiplizieren kann.
}

\begin{myproof}
  Seien $a,b\in\bbN$, wobei $a$ höchstens $k$
  und $b$ höchstens $l$ Binärstellen besitzt.

  Wir zerlegen $b$ in Blöcke mit jeweils
  höchstens $k$ Binärstellen. Setze
  \[
    m:=\left\lceil\frac lk\right\rceil.
  \]
  Dann existieren $b_0,\dots,b_{m-1}$ mit
  $0\le b_i<2^k$, sodass
  \[
    b=\sum_{i=0}^{m-1}b_i2^{ik}.
  \]

  Durch Multiplikation mit $a$ folgt
  \[
    ab=\sum_{i=0}^{m-1}(ab_i)2^{ik}.
  \]

  Jedes Teilprodukt $ab_i$ kann mit Karatsubas
  Algorithmus in
  \[
    O(k^{\log_2 3})
  \]
  Zeit berechnet werden, da beide Faktoren
  höchstens $k$ Binärstellen besitzen.

  Für die insgesamt $m$ Teilprodukte ergibt sich
  somit eine Laufzeit von
  \[
    O\left(mk^{\log_2 3}\right)
    =
    O\left(\frac lk\,k^{\log_2 3}\right)
    =
    O\left(lk^{\log_2 3-1}\right).
  \]

  Es bleibt zu zeigen, dass das Zusammensetzen
  der Teilprodukte nicht mehr Zeit benötigt.

  Jedes Teilprodukt $ab_i$ besitzt höchstens
  $2k$ Binärstellen.

  Wir teilen die Summe in gerade und ungerade
  Blockindizes auf:
  \begin{align*}
    S_0&:=\sum_{\substack{0\le i<m\\i\text{ gerade}}}
      (ab_i)2^{ik},\\
    S_1&:=\sum_{\substack{0\le i<m\\i\text{ ungerade}}}
      (ab_i)2^{ik}.
  \end{align*}

  Innerhalb jeder Summe überlappen die
  Binärblöcke nicht, da aufeinanderfolgende
  Indizes derselben Parität einen Abstand
  von mindestens $2$ haben.

  Somit können $S_0$ und $S_1$ durch
  Aneinanderfügen der Blöcke in $O(l)$
  Zeit gebildet werden.

  Die abschließende Addition
  \[
    ab=S_0+S_1
  \]
  benötigt ebenfalls $O(l)$ Zeit.

  Insgesamt gilt daher
  \begin{align*}
    T(k,l)
    &\in O\left(
      \left\lceil\frac lk\right\rceil
      k^{\log_2 3}+l
    \right)\\
    &=O\left(lk^{\log_2 3-1}\right),
  \end{align*}
  da $k\le l$ und $\log_2 3>1$.

  Damit ist die Behauptung bewiesen.
\end{myproof}


\newpage

\qs{}{
  Zeigen Sie, dass es für alle natürlichen Zahlen
  $m$ und $n$ ganze Zahlen $s$ und $t$ gibt,
  für die
  \[
    \operatorname{ggT}(m,n)=s\cdot m+t\cdot n
  \]
  gilt.
}

\sol
\begin{myproof}
  Wir verwenden den erweiterten euklidischen Algorithmus.

  Zu Beginn gilt
  \[
    m=1\cdot m+0\cdot n,
    \qquad
    n=0\cdot m+1\cdot n.
  \]

  Wir zeigen, dass jeder Rest des euklidischen
  Algorithmus eine ganzzahlige Linearkombination
  von $m$ und $n$ ist.

  Seien dazu zwei aufeinanderfolgende Reste
  \[
    a=s_1m+t_1n,
    \qquad
    b=s_2m+t_2n
  \]
  mit $s_1,s_2,t_1,t_2\in\bbZ$.

  Nach Division mit Rest gilt für ein $q\in\bbZ$:
  \[
    r=a-qb.
  \]

  Durch Einsetzen erhalten wir
  \begin{align*}
    r
    &=s_1m+t_1n-q(s_2m+t_2n)\\
    &=(s_1-qs_2)m+(t_1-qt_2)n.
  \end{align*}

  Da $s_1-qs_2,t_1-qt_2\in\bbZ$,
  besitzt auch $r$ die gewünschte Darstellung.

  Der letzte von null verschiedene Rest
  des euklidischen Algorithmus ist
  $\operatorname{ggT}(m,n)$.

  Somit existieren $s,t\in\bbZ$ mit
  \[
    \boxed{\operatorname{ggT}(m,n)=sm+tn.}
  \]
\end{myproof}
\newpage

\qs{}{
  Beweisen Sie die Korrektheit der folgenden Funktion
  zur Berechnung des größten gemeinsamen Teilers
  zweier nichtnegativer ganzer Zahlen.

  Diese Funktion kommt ohne allgemeine Division
  mit Rest aus. Sie benutzt neben Vergleichen und
  Subtraktionen nur Multiplikationen und Divisionen
  mit $2$ sowie Tests auf gerade bzw. ungerade Zahlen,
  die in Binärdarstellung besonders einfach und schnell
  realisierbar sind.

  Welche Laufzeit hat diese Funktion für zwei Zahlen
  mit höchstens $l$ Binärstellen?

  Vergleichen Sie diese mit der Laufzeitschranke
  für den euklidischen Algorithmus.

  Warum werden in Zeile 10 die Argumente der Funktion
  \texttt{gcd2} vertauscht und wird nicht
  \texttt{gcd2(a / 2, b)} aufgerufen?
}


\begin{codeblock}{C++}
// compute greatest common divisor
// using the binary Euclidean algorithm

int gcd2(int a, int b) {
    if (b == 0) {
        return a;
    }
    if (a % 2 == 0 and b % 2 == 0) {
        // both even
        return 2 * gcd2(a / 2, b / 2);
    }
    if (a % 2 == 0) {
        // only a even
        return gcd2(b, a / 2);
    }
    if (b % 2 == 0) {
        // only b even
        return gcd2(a, b / 2);
    }
    if (a < b) {
        // both odd
        return gcd2(a, (b - a) / 2);
    }
    return gcd2(b, (a - b) / 2);
}
\end{codeblock}

\sol

\begin{myproof}
  Wir zeigen zunächst die Korrektheit und Terminierung
  des Algorithmus und untersuchen anschließend
  seine Laufzeit.

  \medskip
  \textbf{1. Korrektheit:}

  Wir zeigen, dass jeder rekursive Aufruf
  den größten gemeinsamen Teiler erhält,
  wobei im ersten Fall ein Faktor $2$
  berücksichtigt werden muss.

  \textbf{Fall 1:} $b=0$.

  Dann gilt
  \[
    \operatorname{ggT}(a,0)=a.
  \]
  Der Algorithmus gibt also den korrekten Wert zurück.

  \textbf{Fall 2:} $a$ und $b$ sind gerade.

  Dann gilt
  \[
    \operatorname{ggT}(a,b)
    =2\operatorname{ggT}\left(
      \frac a2,\frac b2
    \right).
  \]
  Der rekursive Aufruf liefert somit nach
  Multiplikation mit $2$ den gesuchten ggT.

  \textbf{Fall 3:} Genau eine Zahl ist gerade.

  Ist $a$ gerade und $b$ ungerade, so gilt
  \[
    \operatorname{ggT}(a,b)
    =\operatorname{ggT}\left(
      b,\frac a2
    \right).
  \]
  Denn jeder gemeinsame Teiler muss ungerade
  sein, sodass der Faktor $2$ entfernt
  werden kann.

  Ist dagegen $b$ gerade und $a$ ungerade,
  gilt entsprechend
  \[
    \operatorname{ggT}(a,b)
    =\operatorname{ggT}\left(
      a,\frac b2
    \right).
  \]

  \textbf{Fall 4:} $a$ und $b$ sind ungerade.

  Für $a<b$ gilt
  \begin{align*}
    \operatorname{ggT}(a,b)
    &=\operatorname{ggT}(a,b-a)\\
    &=\operatorname{ggT}\left(
      a,\frac{b-a}{2}
    \right).
  \end{align*}

  Analog gilt für $a\ge b$:
  \[
    \operatorname{ggT}(a,b)
    =\operatorname{ggT}\left(
      b,\frac{a-b}{2}
    \right).
  \]

  Somit liefern alle rekursiven Aufrufe
  den korrekten größten gemeinsamen Teiler.

  \medskip
  \textbf{2. Terminierung:}

  Bei jedem rekursiven Aufruf mit
  $a,b>0$ wird mindestens eine der Zahlen
  echt kleiner, während die andere
  nicht größer wird.

  Insbesondere sinkt die Summe $a+b$
  bei jedem solchen Schritt strikt.

  Da $a+b\in\bbN$ ist, kann dies
  nicht unendlich oft geschehen.

  Falls $a=0$ und $b>0$, so wird
  entweder $b$ halbiert oder direkt
  der Aufruf $\operatorname{gcd2}(b,0)$
  erreicht.

  Damit wird nach endlich vielen Schritten
  der Abbruchfall $b=0$ erreicht.
  Der Algorithmus terminiert also.

  Zusammen mit der Korrektheit der
  rekursiven Schritte folgt, dass
  die Funktion den ggT berechnet.

  \medskip
  \textbf{3. Laufzeit:}

  Seien $a$ und $b$ Zahlen mit höchstens
  $l$ Binärstellen.

  Wir betrachten die Summe der Bitlängen
  \[
    L(a,b):=\ell(a)+\ell(b),
  \]
  wobei $\ell(0):=0$ gilt.

  Bei jedem rekursiven Aufruf mit
  $a,b>0$ wird mindestens eine positive
  Zahl durch eine Zahl ersetzt, die
  höchstens halb so groß ist.

  Daher sinkt $L(a,b)$ bei jedem
  solchen Aufruf um mindestens $1$.

  Anfangs gilt
  \[
    L(a,b)\le2l.
  \]

  Folglich gibt es höchstens $O(l)$
  rekursive Aufrufe.

  Jeder Aufruf benötigt für Vergleiche,
  Subtraktionen und Halbierungen höchstens
  $O(l)$ Bitoperationen.

  Somit beträgt die Gesamtlaufzeit
  \[
    \boxed{T(l)\in O(l^2).}
  \]

  Der gewöhnliche euklidische Algorithmus
  besitzt mit schriftlicher Division die
  Laufzeitschranke $O(l^3)$.

  Der binäre euklidische Algorithmus
  hat somit eine bessere asymptotische
  obere Schranke.

  \medskip
  \textbf{4. Vertauschung der Argumente:}

  In Zeile 10 wird
  \[
    \operatorname{gcd2}\left(b,\frac a2\right)
  \]
  statt
  \[
    \operatorname{gcd2}\left(\frac a2,b\right)
  \]
  aufgerufen.

  Die Vertauschung ist für die
  Terminierung bei $a=0$ notwendig.

  Beispielsweise würde ohne Vertauschung
  der Aufruf
  \[
    \operatorname{gcd2}(0,3)
  \]
  immer wieder sich selbst aufrufen.

  Mit Vertauschung entsteht dagegen
  \[
    \operatorname{gcd2}(3,0),
  \]
  sodass der Abbruchfall erreicht wird.
\end{myproof}

\newpage

\qs{}{
  Beweisen Sie die folgenden Aussagen
  über die Fibonacci-Zahlen:
  \begin{enumerate}
    \item Die \textbf{Cassini-Identität}:
      \[
        F_{n+1}F_{n-1}-F_n^2=(-1)^n,
        \qquad n\ge1.
      \]

    \item Für $k\ge1$ und $n\ge0$ gilt:
      \[
        F_{n+k}
        =F_kF_{n+1}+F_{k-1}F_n.
      \]

    \item Für $k\ge1$ ist $F_{kn}$
      ein Vielfaches von $F_n$.

    \item Für $m,n\ge1$ gilt:
      \[
        \operatorname{ggT}(F_m,F_n)
        =F_{\operatorname{ggT}(m,n)}.
      \]
  \end{enumerate}

  \textbf{Hinweis zu (d):}

  Beweisen Sie zunächst unter Verwendung von
  (b) und (c), dass
  \[
    \operatorname{ggT}(F_m,F_n)
    =\operatorname{ggT}(F_n,F_r)
  \]
  gilt, falls $m$ mit Rest $r$ durch $n$
  teilbar ist.

  Beachten Sie außerdem, dass je zwei
  aufeinanderfolgende Fibonacci-Zahlen
  teilerfremd sind.
}

\newpage

\qs{}{
  Ergänzen Sie in der \texttt{add}-Funktion
  der Klasse \texttt{Fraction} einen Test,
  der vor jeder Rechenoperation prüft,
  ob ein Überlauf auftreten würde, und
  in diesem Fall mittels \texttt{throw}
  eine Ausnahme auslöst.

  Beachten Sie, dass ein Überlauf sowohl
  in einer der drei Multiplikationen
  als auch bei der Addition auftreten kann.

  Nennen Sie die so modifizierte
  \texttt{add}-Funktion \texttt{add\_safe}.
}

\sol

\end{document}
